\newpage
\thispagestyle{plain}

% ── Résumé en Français ────────────────────────────────────────────────────────
\section*{Résumé}
\addcontentsline{toc}{chapter}{Résumés / Abstract}

Dans un contexte mondial marqué par la recrudescence des cyberattaques étatiques et le durcissement du cadre réglementaire national impulsé par le \textbf{Décret-loi n° 2023-17}, la Tunisie fait face au défi stratégique de la souveraineté de son espace numérique. Ce projet de fin d'année (PFA) présente la conception et l'implémentation de \textbf{CyberScore TN}, une plateforme nationale et souveraine d'évaluation continue de la posture de cybersécurité des 120+ institutions gouvernementales et souveraines du domaine public (\texttt{.gov.tn}, \texttt{.tn} régaliens et \texttt{.com.tn}).

Reposant sur un moteur asynchrone Python (FastAPI / AsyncIO) et une interface Next.js 16 trilingue (Arabe RTL, Français, Anglais), CyberScore TN orchestre \textbf{16 piliers d'audit passifs non-intrusifs} (SSL/TLS TunTrust, En-têtes HTTP, DNSSEC, DMARC, CSP, Surface d'attaque EASM, Supply Chain, Threat Intelligence). Les vulnérabilités sont automatiquement corrélées avec les référentiels \textbf{OWASP Top 10}, \textbf{NIST CSF v2} et \textbf{ISO 27001}, alimentant un algorithme de scoring dual (Note A--F et label 3 étoiles). La plateforme innove par l'intégration d'un \textbf{LLM local souverain (Qwen2.5-Coder:7b sous Ollama)} sans fuite cloud, d'un avatar 3D interactif (\textit{Amina}), de projections prédictives de risque à +30 jours, d'un centre de gouvernance RBAC cloisonné et d'un générateur de rapports officiels vectoriels certifiés par empreinte cryptographique SHA-256.

\vspace{0.3cm}
\noindent\textbf{Mots-clés :} Cybersécurité Souveraine, Audit Passif, Décret 2023-17, Scoring Prédictif, IA Locale Ollama, Next.js 16, FastAPI, ANCS, TEK-UP University, Tunisie, PFA.

\vspace{0.8cm}
\hrule
\vspace{0.8cm}

% ── Abstract en Anglais ───────────────────────────────────────────────────────
\section*{Abstract}

In a global landscape characterized by escalating state-sponsored cyber threats and tighter national compliance under \textbf{Decree-Law No. 2023-17}, Tunisia must safeguard the sovereignty and digital resilience of its public institutions. This engineering project (PFA) introduces \textbf{CyberScore TN}, a sovereign national platform designed for automated, continuous, and non-intrusive cybersecurity posture assessment across all 120+ public administration and sovereign domains (\texttt{.gov.tn}, sovereign \texttt{.tn}, and \texttt{.com.tn}).

Built with an asynchronous FastAPI backend and a Next.js 16 trilingual frontend (Arabic RTL, French, English), CyberScore TN executes \textbf{16 passive inspection pillars} (TunTrust SSL/TLS, Security Headers, DNSSEC, DMARC, CSP, EASM attack surface, Supply Chain, Threat Intelligence). Detected issues are correlated with \textbf{OWASP Top 10}, \textbf{NIST CSF v2}, and \textbf{ISO 27001}, powering a dual scoring algorithm (Letter grades A--F + 3-star seal). Key innovations include a \textbf{100\% on-premise sovereign LLM (Qwen2.5-Coder:7b via Ollama)} preventing foreign data leakage, an interactive 3D virtual assistant (\textit{Amina}), 30-day predictive risk modeling, role-based governance (RBAC with 403 enforcement), and automated generation of cryptographic SHA-256 certified audit PDF reports.

\vspace{0.3cm}
\noindent\textbf{Keywords :} Sovereign Cybersecurity, Passive Scanning, Predictive Scoring, Local LLM, Non-Intrusive Audit, Next.js 16, FastAPI, ANCS, TEK-UP University, Gov.tn, PFA.
